\section{The centralizer and the full automorphism group}

Section~5 established that the graph

\[
X
\]

admits a homogeneous-space realization as the connected quartic orbital
graph arising from the action of

\[
S_{5}
\]

on

\[
S_{5}/V_{4,\mathrm{even}}.
\]

This realization exposes symmetries that are invisible from the canonical
covering construction alone.

In this section we determine the full automorphism group of \(X\).


\subsection{The normalizer of \(V_{4,\mathrm{even}}\)}

% STATUS: Conceptual

Let
\[
H
=
V_{4,\mathrm{even}}
=
\{
1,
(12)(34),
(13)(24),
(14)(23)
\}.
\]

\begin{proposition}
\label{prop:normalizer-quotient}
The normalizer of \(H\) in \(S_5\) is the point stabilizer of \(5\).
Consequently,
\[
N_{S_5}(H)\cong S_4
\]
and
\[
N_{S_5}(H)/H\cong S_3.
\]
\end{proposition}

\begin{proof}
The common fixed-point set of the three nonidentity elements of \(H\) is
the singleton \(\{5\}\). Every element normalizing \(H\) must preserve
this set, so
\[
N_{S_5}(H)
\leq
\operatorname{Stab}_{S_5}(5)
\cong
S_4.
\]

Conversely, the point stabilizer of \(5\) permutes the three double
transpositions in \(H\), and therefore normalizes \(H\). Hence
\[
N_{S_5}(H)
=
\operatorname{Stab}_{S_5}(5)
\cong
S_4.
\]

Conjugation on the three nonidentity elements of \(H\) gives a surjective
homomorphism
\[
N_{S_5}(H)\longrightarrow S_3.
\]
Its kernel is the centralizer of \(H\) in \(S_4\), which is \(H\) itself.
Therefore
\[
N_{S_5}(H)/H
\cong
S_3.
\]
\end{proof}

\subsection{The centralizer of the coset action}

% STATUS: Conceptual

Let
\[
\Omega
=
S_5/H
\]
with \(H=V_{4,\mathrm{even}}\), and let \(S_5\) act by left
multiplication.

\begin{corollary}
\label{cor:coset-centralizer}
The centralizer of this permutation representation is
\[
C_{\Sym(\Omega)}(S_5)
\cong
S_3.
\]
\end{corollary}

\begin{proof}
By Lemma~\ref{lem:coset-centralizer},
\[
C_{\Sym(\Omega)}(S_5)
\cong
N_{S_5}(H)/H.
\]
Proposition~\ref{prop:normalizer-quotient} identifies the quotient with
\(S_3\).
\end{proof}

\subsection{Fixation of the distinguished orbital}

% STATUS: Conceptual deduction from certified computation
% DEPENDS: Certified quartic orbital census in Section 5
% LITERATURE: Cameron2003CoherentConfigurations
%             Eberhard2023HammingSandwiches

The centralizer may permute the orbital relations of the Schurian coherent
configuration; it need not fix every relation individually. The certified
uniqueness result of Section~5 forces fixation in the present case.

\begin{theorem}
\label{thm:centralizer-fixes-orbital}
The centralizer of the \(S_5\)-action on
\[
\Omega=S_5/V_{4,\mathrm{even}}
\]
acts by automorphisms of the distinguished connected quartic orbital
graph. Consequently,
\[
C_{\Sym(\Omega)}(S_5)
\leq
\Aut(X).
\]
\end{theorem}

\begin{proof}
The centralizer is contained in the normalizer of the permutation group.
By Lemma~\ref{lem:normalizer-permutes-orbitals}, it permutes the
self-paired orbital relations and carries each orbital graph to an
isomorphic orbital graph.

The certified census in Section~5 contains exactly four quartic
self-paired orbital graphs. Exactly one is connected; the other three each
have two components of order fifteen. Valency and connectedness are
preserved under graph isomorphism, so the unique connected quartic
relation is fixed by every element of the centralizer.

By Theorem~\ref{thm:coincidence}, this distinguished orbital graph
is \(X\). Hence every element of the centralizer preserves the edge
relation of \(X\).
\end{proof}

Combining Theorem~\ref{thm:centralizer-fixes-orbital} with
Corollary~\ref{cor:coset-centralizer} gives
\[
S_3
\cong
C_{\Sym(\Omega)}(S_5)
\leq
\Aut(X).
\]

\subsection{The commuting product}

% STATUS: Conceptual deduction

Under the homogeneous-space identification of Section~5, let

\[
K\cong S_{5}
\]

denote the subgroup induced by left multiplication on

\[
\Omega=S_{5}/V_{4,\mathrm{even}},
\]

and let

\[
C=C_{\Sym(\Omega)}(K).
\]

The orbital construction gives

\[
K\leq\Aut(X),
\]

while the fixation theorem of the preceding subsection gives

\[
C\leq\Aut(X).
\]

By Corollary~\ref{cor:coset-centralizer},

\[
C\cong S_{3}.
\]

\begin{proposition}
\label{prop:commuting-product}

The subgroups \(K\) and \(C\) generate an internal direct product

\[
KC
\cong
S_{5}\times S_{3}.
\]

In particular,

\[
|KC|=720.
\]

\end{proposition}

\begin{proof}

By definition of the centralizer, every element of \(C\) commutes with
every element of \(K\).

It remains to determine the intersection. If

\[
x\in K\cap C,
\]

then \(x\) belongs to \(K\) and centralizes all of \(K\). Hence

\[
x\in Z(K).
\]

Since

\[
K\cong S_{5}
\]

and \(S_{5}\) has trivial center,

\[
K\cap C=1.
\]

Therefore the commuting product is internal and direct:

\[
KC
\cong
K\times C
\cong
S_{5}\times S_{3}.
\]

Its order is

\[
|KC|
=
|S_{5}|\cdot|S_{3}|
=
120\cdot6
=
720.
\]
\end{proof}

\subsection{The full automorphism group}

% STATUS: Certified computational input followed by conceptual deduction
% SOURCE: all_one_group_anatomy_audit_028o.json
% COMMIT: 1604b16

The remaining computational input is an exact enumeration of the full
automorphism group.

\begin{theorem}
\label{thm:full-automorphism-group}

The full automorphism group of \(X\) is

\[
\Aut(X)
\cong
S_{5}\times S_{3}.
\]

\end{theorem}

\begin{proof}

The certified automorphism computation described in Section~9 gives

\[
|\Aut(X)|=720.
\]

By Proposition~\ref{prop:commuting-product}, the group \(\Aut(X)\)
already contains the subgroup

\[
KC
\cong
S_{5}\times S_{3}
\]

of order \(720\).

Consequently,

\[
\Aut(X)=KC,
\]

and therefore

\[
\Aut(X)
\cong
S_{5}\times S_{3}.
\]
\end{proof}

The computation is used only to exclude automorphisms outside the
conceptually constructed subgroup \(K\times C\). The complementary
\(S_{3}\) itself is forced by the homogeneous-space realization and the
uniqueness of the connected quartic orbital relation.

\subsection{Group-theoretic consequences}

% STATUS: Conceptual deduction from Theorem 6.5

The direct-product description determines the elementary anatomy of the
full automorphism group.

\begin{corollary}
\label{cor:group-anatomy}

Let

\[
A=\Aut(X).
\]

Then

\[
Z(A)=1,
\]

\[
A'
\cong
A_{5}\times C_{3},
\]

and

\[
A^{\mathrm{ab}}
\cong
C_{2}\times C_{2}.
\]

In particular,

\[
|A'|=180
\]

and

\[
[A:A']=4.
\]

\end{corollary}

\begin{proof}

By Theorem~\ref{thm:full-automorphism-group},

\[
A
\cong
S_{5}\times S_{3}.
\]

The center of a direct product is the product of the centers. Since both
\(S_{5}\) and \(S_{3}\) have trivial center,

\[
Z(A)
=
Z(S_{5})\times Z(S_{3})
=
1.
\]

The derived subgroup of a direct product is the product of the derived
subgroups. Using

\[
S_{5}'=A_{5}
\]

and

\[
S_{3}'=A_{3}\cong C_{3},
\]

we obtain

\[
A'
\cong
A_{5}\times C_{3}.
\]

Finally,

\[
S_{5}^{\mathrm{ab}}\cong C_{2}
\]

and

\[
S_{3}^{\mathrm{ab}}\cong C_{2},
\]

so

\[
A^{\mathrm{ab}}
\cong
C_{2}\times C_{2}.
\]

The order and index statements follow immediately.
\end{proof}

The full group is therefore centerless, while its derived subgroup retains
the alternating symmetry of the Petersen structure together with the
three-cycle subgroup of the complementary \(S_{3}\).

The next section interprets the three involutions of that complementary
factor as an automorphism orbit of regular quotient frames.
